\section*{Efectos Estocásticos y Protección Radiológica}

Los efectos estocásticos son aquellos que aparentemente no tienen una dosis umbral a partir de la cual se manifiestan, la probabilidad de que ocurra el efecto independientemente de la severidad, se considera en función de la dosis, además debe tenerse en cuenta que la mayoría de estos efectos existen normalmente en poblaciones no irradiadas. Estos efectos se pueden dividir en cuatro grupos:
\begin{itemize}
    \item Inducción de leucemia.
    \item Inducción de otros cánceres.
    \item Enfermedades hereditarias.
    \item Acortamiento de la vida.
\end{itemize}

El propósito de la protección radiológica es evitar los efectos determinísticos fijando límites de dosis por debajo de sus umbrales.

\section*{Formas de Exposición a la Radiación}

Dentro de las formas de exposición a la Radiación, se encuentra la irradiación externa la que recibe el organismo desde fuentes exteriores al mismo, tenuto por ser consideran fáciles de detectar mas peligrosas, tanto por ser los más difíciles de detectar, como por su efecto sobre los tejidos, ya que al interaccionar con los átomos de hidrógeno producen protones cuyo elevado poder de ionización causa un daño muy localizado.

La irradiación interna se presenta cuando la fuente se encuentra dentro del organismo. Los materiales radiactivos pueden ingresar al cuerpo por ingestión, por inhalación, por absorción a través de la piel o bien a la sangre por algún corte o herida. El daño depende del tipo de radiación emitida en su energía, la vía media de radionúclido, su permanencia en el organismo y el órgano en que se localice. En este caso los emisores alfa y beta son los más peligrosos.

A continuación se anexan cuadros que resumen los efectos de la irradiación al organismo.
